%% secoes/introducao/04-apgar.tex — incluído por main.tex

%\subsection{Apgar baixo no quinto minuto}\label{subsec:intro-apgar}

Por fim, o índice de Apgar, proposto por Virginia Apgar em 1953, avalia a condição do recém-nascido logo após o parto a partir de cinco sinais: frequência cardíaca, esforço respiratório, tônus muscular, irritabilidade reflexa e cor da pele, cada um pontuado de 0 a 2 \cite{apgar1953}. A avaliação é feita no primeiro e no quinto minuto de vida, e o índice inferior a 7 no quinto minuto é usado como indicador de comprometimento das condições de nascimento e da necessidade de cuidados especializados \cite{magalhaes2023}.
No Brasil, a proporção de recém-nascidos com Apgar menor que 7 no quinto minuto caiu de 2,1\% em 1999 para 0,9\% em 2018--2019 \cite{magalhaes2023}. 
Entre os quatro desfechos deste trabalho, o Apgar baixo é o menos investigado sob a perspectiva social (Seção~\ref{sec:relacionados}).

%Prematuridade, baixo peso ao nascer e anomalia congênita são os fatores de maior impacto, mas também se associam ao desfecho características maternas e da assistência, como idade a partir de 35 anos, escolaridade inferior a oito anos, ausência de companheiro e menos de seis consultas de pré-natal \cite{magalhaes2023}. Apesar disso, entre os quatro desfechos deste trabalho, o Apgar baixo é o menos investigado sob a perspectiva social (Seção~\ref{sec:relacionados}).
